% Exemplo do perfil classico da classe gmap-ufrgs.
% Compile: latexmk -pdf exemplo-classico.tex   (dados ficticios)
% No perfil classico a capa e o sumario saem sozinhos no \begin{document}:
% o texto comeca direto no capitulo. As listas de simbolos e de siglas e o
% glossario saem sozinhos no fim do documento.
\documentclass[classico]{gmap-ufrgs}

\projeto{Projeto Exemplo}
\numerodoc{GMAp-RT-000/2026}
\revisao{0}
\titulo{Exemplo do Perfil Cl\'assico}
\local{Porto Alegre}
\addautor{Beltrana de Tal}{Eng. Mec\^anica}{CREA-RS 000001}
\equipe{Prof. Fulano de Tal\\ Beltrana de Tal}

%--------------------------------------------------------------------
% Siglas, abreviaturas, simbolos e glossario (glossaries-extra).
% So entram nas listas os itens usados no texto. Podem ficar num arquivo
% a parte, lido no preambulo com \input{siglas}.
%--------------------------------------------------------------------
% Siglas: 1.a vez "Expansao (SIGLA)", depois so "SIGLA"
\sigla{mef}{MEF}{M\'etodo dos Elementos Finitos}
\sigla[user1={Applied Mechanics Group}]{gmap}{GMAp}{Grupo de Mec\^anica Aplicada}
\newacronym{abnt}{ABNT}{Associa\c{c}\~ao Brasileira de Normas T\'ecnicas}
% Abreviaturas: sempre so a forma curta
\abreviatura{etal}{et al.}{et alii}
% Simbolos (lista no fim do documento, na ordem em que aparecem no
% texto). Use \ensuremath,
% nao $...$: assim \gls{sigma} funciona no texto e dentro de equacoes
\simbolo{sigma}{\ensuremath{\sigma}}{Tens\~ao normal}
\glsxtrnewsymbol[description={M\'odulo de elasticidade}]{E}{\ensuremath{E}}
% Glossario (no fim do documento)
\termo{malha}{malha}{Divis\~ao do dom\'inio em elementos.}
\newglossaryentry{no}{name={n\'o},description={Ponto que liga elementos da malha.},sort={no}}

\begin{document}
\chapter{Introdu\c{c}\~ao}
Este relat\'orio do \gls{gmap} usa o \gls{mef}, conforme a \gls{abnt}. Na
segunda vez, s\'o a sigla: \gls{mef}. Segundo Silva \gls{etal}, a
\gls{malha} tem tens\~ao \gls{sigma} e m\'odulo \gls{E} em cada \gls{no}.

\figura[fonte={Elaborada pelos autores.}]{example-image-a}{Figura de exemplo.}

% nofim: a assinatura sai no fim de tudo (depois das listas);
% novapagina: numa pagina so dela
\blocoassinatura[nofim,novapagina]{Prof. Fulano de Tal\\ Coordenador do projeto\\
  Grupo de Mec\^anica Aplicada}
\end{document}
